\chapter{Dynamic and Kinetic Geometry}
\label{ch:dynamic_kinetic_geometry}

\section{Dynamic Geometric Problems}
\label{sec:dynamic_intro}

Every data structure in the preceding chapters was built for a \emph{static}
input: a set of points, segments, or polygons fixed once and for all. Real
applications rarely respect that assumption: terrain vertices change as new
survey data arrives, database entries are inserted and deleted, and the
robots in a warehouse move continuously. This chapter studies geometric
algorithms whose input \emph{changes}, and it distinguishes three regimes:

\begin{description}
  \item[Static] The input is fixed; queries are answered against a structure
    built once (Chapter~\ref{ch:range_searching}).
  \item[Dynamic] The input changes by discrete insertions and deletions
    interleaved with queries (Sections~\ref{sec:dynamic_point_sets}--\ref{sec:dynamic_proximity}).
  \item[Kinetic] The input changes \emph{continuously} along known
    trajectories, and the answer is a function of time maintained by
    \emph{kinetic data structures} (Sections~\ref{sec:dynamic_kds}--\ref{sec:dynamic_kinetic_closest}).
\end{description}

A dynamic structure is judged by its \textbf{update time}\index{update time},
\textbf{query time}, and \textbf{storage}. A kinetic structure is judged by
the additional quantities of Section~\ref{sec:dynamic_kds}: how many
\emph{certificates} it holds, how fast it repairs a \emph{certificate
failure}, and how the number of failures compares with the number of times
the answer changes. The unifying theme is \textbf{coherence}\index{coherence}:
between updates, or between small advances of time, the combinatorial answer
barely changes, and a good structure changes only what it must. The dynamic
results rest on the transform-and-conquer machinery of Overmars
\cite{overmars1983}; the kinetic results rest on the kinetic data structure
framework of Basch, Guibas, and Hershberger \cite{basch1999,guibas1998}.

\index{dynamic data structure}\index{kinetic data structure}

\section{Dynamic Point Sets}
\label{sec:dynamic_point_sets}

The simplest dynamic geometric object is a set of points supporting
\textbf{insert} and \textbf{delete} operations interleaved with queries.
Most point-based structures --- $k$-d trees, range trees, Voronoi diagrams
--- have dynamic versions, and the classical framework for deriving them is
the \textbf{transform-and-conquer} method of Overmars \cite{overmars1983}.

\begin{definition}[Dynamic Point Structure]\label{def:dynamic_structure}
  A \textbf{dynamic point structure} maintains a set $P \subseteq
  \mathbb{R}^d$ under $\textsc{Insert}(p)$, $\textsc{Delete}(p)$, and a class
  of queries. It is \textbf{semidynamic}\index{semidynamic} if it supports
  only insertions and \textbf{fully dynamic} if it supports both.
\end{definition}

The two basic design strategies are:

\begin{itemize}
  \item \textbf{Partial rebuilding}\index{partial rebuilding}. The
    \textbf{logarithmic method} partitions the points into $O(\log n)$
    static structures of geometrically growing size: a new point enters the
    smallest structure, and structures merge when their size doubles. A
    query runs against all levels, so the update cost is an $O(\log n)$
    factor times the static construction cost and the query cost an
    $O(\log n)$ factor times the static query cost, both amortized.
    Section~\ref{sec:rangesearch_dynamic} analyzes this for range searching.
  \item \textbf{Balanced-tree dynamics}. The primary structure is kept
    balanced, and each rotation repairs only local secondary structures. A
    dynamic range tree costs $O(\log^2 n)$ per update in the plane
    (Section~\ref{sec:rangesearch_dynamic}); a dynamic $k$-d tree
    (Section~\ref{sec:kd_tree}) inserts by leaf attachment and deletes by
    marking with periodic subtree rebuilds, at $O(\log n)$ amortized cost.
\end{itemize}

Two auxiliary problems recur. The \textbf{order maintenance}
problem\index{order maintenance problem} --- maintain a sorted list so that
the relative order of two elements is tested in $O(1)$ --- is solved by the
order-maintenance structure of Dietz and Sleator \cite{dietzsleator1987},
which gives each element a monotone label from a dense order. The second is
\textbf{dynamic point location}, answered across updates at expected
$O(\log n)$ per insertion and $O(\log^2 n)$ per deletion
(Section~\ref{sec:pointloc_dynamic}, \cite{overmars1983}). Textbook
treatments are \cite{overmars1983,berg2008} and the handbook chapters of
\cite{mehta2005}; the priority search tree of
Section~\ref{sec:rangesearch_priority_search} is the classic linear-storage
structure for dynamic 3-sided queries.

\index{order maintenance problem}

\section{Dynamic Convex Hulls}
\label{sec:dynamic_convex_hulls}

\subsection{1. Overview}

The dynamic convex hull problem maintains $\operatorname{conv}(P)$ under
insertions and deletions while answering queries about the hull. The two
classical milestones are the $O(\log^2 n)$-per-update structure of Overmars
and van Leeuwen \cite{overmars1981} and the near-optimal structure of Brodal
and Jacob \cite{brodal2002}, whose total cost over any sequence of $n$
updates is $O(n\log n)$. This section develops the Overmars--van Leeuwen
approach and states the Brodal--Jacob improvement.

\subsection{2. Definitions}

\begin{definition}[Dynamic Convex Hull]\label{def:dyn_convex_hull}
  A \textbf{dynamic convex hull} structure maintains $\operatorname{conv}(P)$
  under $\textsc{Insert}(p)$ and $\textsc{Delete}(p)$, together with the
  queries of the following definition.
\end{definition}

\begin{definition}[Hull Queries]\label{def:dyn_extreme_query}
  For a direction $d$, the \textbf{extreme-point query} asks for a vertex of
  $\operatorname{conv}(P)$ maximizing $p \cdot d$; the \textbf{membership
  query} asks whether a given point lies inside the hull. Both are answered
  in $O(\log n)$ by binary search over the cyclic hull order
  (Chapter~\ref{ch:convex_hulls}).
\end{definition}

\subsection{3. Theory}

The Overmars--van Leeuwen structure maintains the upper hull and the lower
hull separately, each as a balanced binary tree over the vertices in
$x$-order. The key idea is the \textbf{bridge}\index{bridge}: if a subtree's
vertices are split at a median $x$-coordinate, the subtree's upper hull is
the upper hull of the left part, one \emph{bridge edge} connecting it to the
upper hull of the right part, and the right part's upper hull. A node stores
its bridge, found in $O(\log n)$ time by a walk that advances a current
vertex alternately on the two child chains.

\begin{theorem}[Overmars--van Leeuwen]\label{thm:dyn_hull_overmars}
  The convex hull of an $n$-point set can be maintained under insertions and
  deletions in $O(\log^2 n)$ worst-case time per update, with extreme-point
  and membership queries in $O(\log n)$ time and $O(n)$ storage.
\end{theorem}

\begin{proof}[Proof sketch]
  An update first rebalances the hull tree. The affected nodes lie on one
  root-to-leaf path, and at each such node the bridge is recomputed from the
  two child hull chains in $O(\log n)$ time; there are $O(\log n)$ nodes,
  giving $O(\log^2 n)$. A direction query walks one path, spending $O(1)$
  per level following bridges, hence $O(\log n)$; the same search supports
  the membership test.
\end{proof}

The structure is the canonical \emph{two-level} dynamic structure, in the
spirit of the multilevel structures of
Section~\ref{sec:rangesearch_multilevel}. Brodal and Jacob \cite{brodal2002}
improved the update bound by a more carefully amortized, partition-based
representation of the hull and its bridges:

\begin{theorem}[Brodal--Jacob]\label{thm:dyn_hull_brodal}
  The convex hull of an $n$-point set can be maintained under insertions and
  deletions in $O(\log n)$ amortized time per update and $O(\log n)$
  worst-case time per query, using $O(n)$ storage. Any sequence of $n$
  updates thus costs $O(n\log n)$ total, which is optimal up to logarithmic
  factors for comparison-based exact structures.
\end{theorem}

\subsection{4. Pseudocode}

\begin{algorithm}[htbp]
\caption{Dynamic Convex Hull Maintenance}
\label{alg:dyn_convex_hull}
\begin{algorithmic}[1]
\Require A dynamic point set $P$ in general position.
\Ensure Insertions, deletions, and direction queries on $\operatorname{conv}(P)$.
\Function{DynHullInsert}{$p$}
  \State insert $p$ into the $x$-ordered trees of the upper and lower hulls
  \For{each node $v$ on the rebalanced paths}
    \State $H[v] \gets \textsc{Bridge}(\text{left}(v), \text{right}(v))$
  \EndFor
\EndFunction
\Function{DynHullDelete}{$p$}
  \State delete $p$ from both hull trees; rebalance and recompute bridges as above
\EndFunction
\Function{ExtremePoint}{$d$}
  \State $v \gets$ root of the upper-hull tree
  \While{$v$ is not a leaf}
    \State follow the bridge stored at $v$ toward the maximum of the dot product with $d$
  \EndWhile
  \State \Return the vertex reached; search the lower hull symmetrically for direction $d$
\EndFunction
\end{algorithmic}
\end{algorithm}

\subsection{5. Parameter Selection}

\begin{itemize}
  \item \textbf{Balancing}: any balanced binary search tree works; the
    requirement is that an update touches only $O(\log n)$ nodes, so bridge
    recomputations stay on one root-to-leaf path.
  \item \textbf{Degenerate points}: vertical alignments, duplicate
    $x$-coordinates, and collinear hull vertices are handled by a symbolic
    perturbation (Section~\ref{sec:intro_symbolic_perturbation}) or by
    storing, at each $x$, the interval of hull vertices at that $x$.
\end{itemize}

\subsection{6. Complexity}

\begin{itemize}
  \item \textbf{Update}: $O(\log^2 n)$ worst case (Overmars--van Leeuwen),
    $O(\log n)$ amortized (Brodal--Jacob).
  \item \textbf{Query}: $O(\log n)$; \textbf{storage}: $O(n)$.
  \item \textbf{Total}: $O(n\log^2 n)$ and $O(n\log n)$ over $n$ updates,
    respectively. A comparison-based lower bound of $\Omega(\log n)$ per
    operation makes the Brodal--Jacob bounds optimal up to constants.
\end{itemize}

\subsection{7. Usages}

\begin{itemize}
  \item \textbf{Kinetic convex hulls}: the structure of
    Section~\ref{sec:dynamic_kinetic_hulls} uses a dynamic hull as a
    subroutine when a point enters or leaves the hull.
  \item \textbf{Motion planning}: maintaining obstacle hulls under updates
    supports the queries of Chapter~\ref{ch:motion_planning}.
  \item \textbf{Onion peeling}: repeated extreme-point queries with dynamic
    deletion compute hull layers (Section~\ref{sec:convex_hull_peeling})
    online.
\end{itemize}

\subsection{8. Testing Methods and Edge Cases}

\begin{itemize}
  \item \textbf{Brute-force oracle}: after every update, recompute the hull
    statically and compare vertex sets over random update sequences.
  \item \textbf{Direction queries}: verify \textsc{ExtremePoint}($d$) against
    an $O(n)$ scan over a grid of directions.
  \item \textbf{Degenerate hulls}: collinear points, duplicates, and
    two-point hulls must not break the bridge search.
\end{itemize}

\subsection{9. References}

The $O(\log^2 n)$ structure is due to Overmars and van Leeuwen
\cite{overmars1981}; the dynamic-structure theory is \cite{overmars1983}.
The $O(n\log n)$ total result is \cite{brodal2002}. Textbook treatments
appear in \cite{berg2008,preparata1985} and \cite{mehta2005}.

\section{Dynamic Proximity Structures}
\label{sec:dynamic_proximity}

Proximity queries --- nearest neighbor, closest pair, Voronoi cells,
Delaunay adjacency --- are among the most important dynamic geometric
problems. For \textbf{dynamic nearest neighbor}\index{dynamic nearest
neighbor} the standard options are:

\begin{itemize}
  \item \textbf{Dynamic $k$-d tree}: leaf attachment with periodic rebuilds
    gives $O(\log n)$ amortized updates and the branch-and-bound query of
    Section~\ref{sec:kd_tree} (Chapter~\ref{ch:spatial_structures}).
  \item \textbf{Dynamic Delaunay triangulation}: since the nearest neighbor
    of a point is among the vertices of its Delaunay cell, the dynamic
    Delaunay structure of Section~\ref{sec:dynamic_delaunay}
    (Chapter~\ref{ch:delaunay_triangulations}) answers nearest-neighbor
    queries in expected logarithmic time, with local-flip updates.
  \item \textbf{Reduction to hulls}: in the plane, nearest-neighbor queries
    lift to 3D convex hulls (Section~\ref{sec:delaunay_lifting}); the
    dynamic 3D hull of Chan \cite{chan2006} thereby yields a dynamic planar
    nearest-neighbor structure with polylogarithmic bounds.
\end{itemize}

The \textbf{dynamic closest pair}\index{dynamic closest pair} maintains the
Delaunay triangulation together with a secondary structure reporting its
shortest edge --- the static reduction of Chapter~\ref{ch:proximity}, made
dynamic by updating both levels at each insertion or deletion. Dynamic
\textbf{range searching} is developed in Section~\ref{sec:rangesearch_dynamic}
(Chapter~\ref{ch:range_searching}), where the schemes of
Section~\ref{sec:dynamic_point_sets} yield the $O(\log^2 n)$ and
$O(n^{1-1/d})$ trade-offs. When the motion is continuous, the same problems
reappear in kinetic form and the Delaunay triangulation becomes the central
tool (Sections~\ref{sec:dynamic_kinetic_closest}--\ref{sec:dynamic_kinetic_hulls});
the surveys \cite{guibas1998,mehta2005} collect the state of the art.

\index{dynamic nearest neighbor}\index{dynamic closest pair}

\section{Moving Geometric Objects}
\label{sec:dynamic_moving_objects}

Discrete updates are not the only way an input can change. In simulation,
robotics, and tracking, objects are in \emph{continuous motion}, and the
questions are about time: when do two points cross, when do robots collide,
and how does the nearest-neighbor structure evolve?

\begin{definition}[Trajectory]\label{def:kinetic_trajectory}
  A \textbf{trajectory} for an object is a function $p: [0, T] \to
  \mathbb{R}^d$ giving its position at time $t$. The motion is \textbf{linear}
  (constant velocity) if $p(t) = p_0 + vt$, and \textbf{algebraic} if each
  coordinate is a polynomial (or more generally an algebraic function) of
  $t$. The interval $[0, T]$ is the \textbf{time horizon}.
\end{definition}

\begin{example}\label{ex:kinetic_motivating}
  For $n$ points with linear trajectories, the relative order by
  $x$-coordinate changes only at the finitely many times two points have
  equal $x$-coordinate; between such times the order is fixed. Any quantity
  --- sorted order, convex hull, closest pair --- that depends on the
  configuration through order information is therefore \emph{piecewise
  constant in time}, changing only at a discrete set of \emph{event times}
  determined by the trajectories. This is the combinatorial core of kinetic
  geometry: continuous motion, discrete change.
\end{example}

Two strategies are available. The \emph{sampling} strategy re-evaluates the
answer at every step of a discretized clock; it costs a factor $T/\delta$
in the number of samples and can miss events between samples (the classic
\emph{tunneling} artifact of collision detection). The \emph{event-driven}
strategy computes, for each condition that guarantees the current answer,
the exact future time at which it fails, schedules these times in a priority
queue, and processes the earliest failure; it never misses an event but must
compute failure times exactly, which is feasible only for algebraic
trajectories (Section~\ref{sec:dynamic_certificates}). When trajectories are
known in advance (a ``flight plan''), event-driven structures dominate; when
they are revealed on the fly, a hybrid of sampling and event prediction is
used. The kinetic data structure framework of Section~\ref{sec:dynamic_kds}
formalizes the event-driven strategy.

\index{coherence}\index{event}

\section{Kinetic Data Structures}
\label{sec:dynamic_kds}

The \textbf{kinetic data structure} (KDS)\index{kinetic data structure}
framework of Basch, Guibas, and Hershberger \cite{basch1999,guibas1998}
maintains a combinatorial attribute of moving objects by reducing correctness
to a set of small \emph{certificates} whose validity guarantees the current
answer, and scheduling the times at which the certificates fail.

\begin{definition}[Kinetic Structure]\label{def:kinetic_structure}
  A \textbf{kinetic data structure} for an attribute $A$ of $n$ moving
  objects consists of (i) a combinatorial representation of $A$ correct for
  the current configuration, (ii) a set of \textbf{certificates} --- boolean
  tests on the object positions --- such that $A$ is guaranteed correct as
  long as all certificates hold, and (iii) an \textbf{event queue} storing
  the predicted failure time of each certificate.
\end{definition}

\begin{definition}[Certificate]\label{def:kinetic_certificate}
  A \textbf{certificate} is a predicate $c(t)$ over the trajectories whose
  truth at time $t$ is sufficient for the correctness of the maintained
  attribute at $t$. It \textbf{fails} at the smallest $t' > t$ with $c(t')$
  false; this is its \textbf{failure time}\index{failure time}.
\end{definition}

Certificates are elementary geometric tests: order comparisons,
orientation tests, or in-circle tests. Since trajectories are algebraic,
each certificate function is a polynomial in $t$ of constant degree, and its
failure time is a root of that polynomial (at most degree four for an
in-circle certificate under linear motion; Section~\ref{sec:dynamic_certificates}).
\begin{definition}[KDS Quality Measures]\label{def:kinetic_properties}
  Let $s$ be the number of certificates and $n$ the number of objects.
  \begin{itemize}
    \item \textbf{Compactness}: $s = O(n)$.
    \item \textbf{Locality}\index{locality}: each object participates in
      $O(\log n)$ (ideally $O(1)$) certificates.
    \item \textbf{Responsiveness}: a certificate failure is repaired in
      polylogarithmic time.
    \item \textbf{Efficiency}: the number of certificate failures is within
      a polylogarithmic factor of the number of times the maintained
      attribute changes. The ratio separates \emph{internal} events
      (failures that do not change the answer) from \emph{external} events
      (changes of the answer).
  \end{itemize}
\end{definition}

\begin{example}[Kinetic Tournament]\label{ex:kinetic_tournament}
  Maintain the maximum of $n$ points moving on a line. Arrange the points
  as leaves of a balanced tree; each internal node stores the larger child
  and a certificate $c(t) : \text{left winner} \ge \text{right winner}$. As
  long as every certificate holds, the root stores the maximum. When a
  certificate fails, the loser has overtaken the winner: swap, reschedule
  the node's and its predecessors' certificates, and continue. The structure
  holds $O(n)$ certificates, each point is in $O(\log n)$ of them, and a
  failure is repaired in $O(\log n)$ time --- compact, local, responsive. It
  is the canonical first KDS and the model for the kinetic heaps of
  Section~\ref{sec:dynamic_certificates}.
\end{example}

\begin{remark}
  The KDS framework is a \emph{recipe}, not an algorithm: for a new attribute
  one designs the certificate set (a correctness proof), the repair
  procedure (a local update), and the failure-time computation (a root
  solver) \cite{guibas1998}. The following sections apply the recipe to
  sorting, convex hulls, closest pairs, and collision detection.
\end{remark}

\index{kinetic data structure}\index{certificate}\index{locality}

\section{Certificates and Events}
\label{sec:dynamic_certificates}

A KDS is only as good as its event scheduling. This section develops the
algorithmic support: computing failure times, storing them, and repairing
failures reliably.

\begin{definition}[Kinetic Event]\label{def:kinetic_event}
  A \textbf{kinetic event}\index{kinetic event} at time $t$ is a set of one
  or more certificate failures at $t$. It is \textbf{external} if the
  maintained attribute changes and \textbf{internal} otherwise. The event
  queue stores the predicted failure times of all future certificates.
\end{definition}

\subsection*{Computing failure times}

For linear motion $p(t) = p_0 + vt$, an order or orientation certificate
fails at a root of a degree-one polynomial, and an in-circle certificate
(the determinant of the matrix with rows $(x_i, y_i, x_i^2+y_i^2, 1)$) is a
polynomial of degree at most four in $t$. All such roots are computed in
$O(1)$ arithmetic operations, so every certificate schedules its failure in
constant time. This is why KDS implementations are practical: every
predicted failure time is the root of a low-degree polynomial.

\subsection*{Event scheduling with heaps}

The event queue needs extract-min, insert, and decrease, each in $O(\log
n)$; a binary heap delivers all three (Section~\ref{sec:app_alg_priority_queues},
\cite{cormen2009}) and is the standard choice. The \textbf{kinetic
heap}\index{kinetic heap} is a kinetic tournament over the keys to be
minimized (here, the failure times), used to pick the earliest pending
failure.

\begin{definition}[Kinetic Heap]\label{def:kinetic_heap}
  A \textbf{kinetic heap} is a kinetic tournament whose leaves hold the keys
  (failure times or attribute values). Each internal node stores a winner
  and a certificate that its key does not exceed its sibling's; the root
  holds the minimum at all times.
\end{definition}

\subsection*{Repairing failures}

When a certificate fails at time $t$, the repair has three steps: (i)
\emph{verify} the failure with the exact predicates of
Section~\ref{sec:robust_predicates} --- a floating-point evaluation at $t$
can be inconclusive; (ii) \emph{apply} the local combinatorial update (a
swap, a flip, or a tangent recomputation); and (iii) \emph{reschedule} the
new certificates. Two issues dominate practice. First,
\textbf{simultaneous events}: several certificates may fail together, and
the queue must order them by a fixed secondary key or a symbolic
perturbation must break the tie (Section~\ref{sec:intro_symbolic_perturbation}).
Second, \textbf{numerical drift}: a failure time computed in floating point
is slightly wrong, so the certificate is tested at the scheduled time rather
than trusted; if it has not failed yet, the failure time is recomputed and
the event rescheduled (``re-prediction'') \cite{guibas1998}.

\begin{remark}
  The event machinery is not specific to kinetics: the same pattern --- a
  priority queue of events, local repairs, exact predicates --- drives the
  sweep-line algorithms of Chapter~\ref{ch:segment_intersection}, where the
  event queue is the identical heap structure (Section~\ref{sec:segint_events}).
  The difference is that sweep-line events are \emph{discovered} from
  adjacency, while kinetic events are \emph{predicted} from certificate
  roots.
\end{remark}

\index{kinetic heap}\index{kinetic event}

\section{Kinetic Sorting}
\label{sec:dynamic_kinetic_sorting}

\subsection{1. Overview}

Kinetic sorting maintains the sorted order of $n$ points moving on a line,
as a function of time. It is the simplest nontrivial KDS and the foundation
of the kinetic hull and proximity structures, many of which reduce to
maintaining one or more sorted orders (of points by a coordinate, or of
directions by angle). The order can change $\Theta(n^2)$ times in the worst
case even for constant-velocity motion.

\subsection{2. Definitions}

\begin{definition}[Kinetic Sorted List]\label{def:kinetic_sorted_list}
  A \textbf{kinetic sorted list} (KSL) maintains $n$ points moving on a line
  in sorted order by position. The certificates are the $n-1$ adjacent-pair
  inequalities $c_i(t) : p_{\sigma(i)}(t) \le p_{\sigma(i+1)}(t)$, where
  $\sigma$ is the current order; an order change is exactly a certificate
  failure.
\end{definition}

\begin{definition}[Order Maintenance Label]\label{def:kinetic_order_label}
  An \textbf{order label} for each element is a value monotone in the sorted
  order, comparable in $O(1)$; the order-maintenance structures of
  \cite{dietzsleator1987} provide such labels with $O(1)$ amortized update
  cost.
\end{definition}

\subsection{3. Theory}

The KSL is the textbook KDS \cite{basch1999}: it holds $n-1$ certificates
(compact), each point participates in at most two (locality $O(1)$), and a
failure swaps two points, changing only two adjacencies (responsiveness
$O(1)$). Its limitation is the number of events.

\begin{theorem}[Kinetic Sorted List Bounds]\label{thm:kinetic_sort_events}
  The sorted order of $n$ points on a line can change $\Theta(n^2)$ times in
  the worst case, even for constant-velocity motion. The kinetic sorted list
  maintains the order with $O(n)$ certificates and repairs each order change
  in $O(1)$ time plus $O(\log n)$ for the event queue.
\end{theorem}

\begin{proof}
  For the lower bound, take $n/2$ points far to the left moving right with
  velocity $1$ and $n/2$ points on the right moving with velocity $0$. Every
  moving--stationary pair crosses exactly once, reversing its relative order;
  there are $(n/2)^2 = \Theta(n^2)$ such reversals, each forcing an order
  change. For the upper bound, the certificates are the adjacent pairs; a
  failure is a swap changing two adjacencies, and each new adjacency's
  failure time is the root of a degree-one polynomial, computed in $O(1)$
  (Section~\ref{sec:dynamic_certificates}).
\end{proof}

For random or bounded-velocity motions the number of changes is much
smaller, and the KSL is the standard component of kinetic nearest-neighbor
and range structures.

\subsection{4. Pseudocode}

\begin{algorithm}[htbp]
\caption{Kinetic Sorted List}
\label{alg:kinetic_sort}
\begin{algorithmic}[1]
\Require $n$ points $p_1, \dots, p_n$ on a line with trajectories $p_i(t)$.
\Ensure The sorted order at all times in $[0, T]$.
\Function{KineticSortedList}{$p_1, \dots, p_n$}
  \State sort the points at $t = 0$; let $\sigma$ be the current order
  \State $Q \gets$ empty kinetic heap
  \For{$i = 1$ to $n-1$}
    \State create certificate $c_i$ for the pair $(\sigma(i), \sigma(i+1))$
    \State $Q.\text{insert}(\text{FailureTime}(c_i))$
  \EndFor
  \While{$Q$ is not empty and $Q.\text{min}() \le T$}
    \State $t \gets Q.\text{extract-min}()$; advance the clock to $t$
    \State verify the failing pair exactly; swap $\sigma(i)$ and $\sigma(i+1)$
    \State remove the two certificates of the swapped pair; create the two new ones
    \State push the new certificates' failure times into $Q$
  \EndWhile
  \State \Return $\sigma$ and the recorded change times
\EndFunction
\end{algorithmic}
\end{algorithm}

\subsection{5. Parameter Selection}

\begin{itemize}
  \item \textbf{Trajectory form}: the failure-time computation assumes
    polynomial trajectories; for arbitrary trajectories, roots are found
    numerically and events are re-predicted on verification
    (Section~\ref{sec:dynamic_certificates}).
  \item \textbf{Simultaneous crossings}: order coincident failures by a fixed
    secondary key so the swap sequence is deterministic.
  \item \textbf{Order representation}: use order-maintenance labels
    (\cite{dietzsleator1987}) so other structures test order in $O(1)$.
\end{itemize}

\subsection{6. Complexity}

\begin{itemize}
  \item \textbf{Certificates}: $n-1$ (compact); locality $O(1)$.
  \item \textbf{Responsiveness}: $O(\log n)$ per order change.
  \item \textbf{Total}: $O((n^2 + k)\log n)$ worst case for $k$ order
    changes, with $k \le \binom{n}{2}$; $\Theta(n^2)$ events already occur
    for the two-velocity construction of Theorem~\ref{thm:kinetic_sort_events}.
\end{itemize}

\subsection{7. Usages}

\begin{itemize}
  \item \textbf{Kinetic nearest neighbor}: sorted orders by each coordinate
    give candidate certificates for proximity structures
    (Section~\ref{sec:dynamic_proximity}).
  \item \textbf{Kinetic convex hulls}: the angular order of points around a
    moving vantage point drives the hull structure of
    Section~\ref{sec:dynamic_kinetic_hulls}.
  \item \textbf{Event-driven simulation}: any simulation maintaining a set
    of threshold-crossing times uses a kinetic sorted list over them.
\end{itemize}

\subsection{8. Testing Methods and Edge Cases}

\begin{itemize}
  \item \textbf{Oracle}: sort the positions by brute force at a dense grid of
    times and at every event time; compare with the maintained order.
  \item \textbf{Event validation}: verify with exact predicates that the
    scheduled pair truly crossed and that no pair crossed earlier.
  \item \textbf{Degeneracies}: coincident points and simultaneous crossings
    must produce a consistent final order; floating-point drift must be
    absorbed by the re-prediction rule of Section~\ref{sec:dynamic_certificates}.
\end{itemize}

\subsection{9. References}

The kinetic sorted list is a central example of \cite{basch1999,guibas1998};
the order-maintenance labels are due to \cite{dietzsleator1987}. Its use
inside larger kinetic structures is surveyed in \cite{guibas1998,mehta2005}.

\section{Kinetic Convex Hulls}
\label{sec:dynamic_kinetic_hulls}

\subsection{1. Overview}

The kinetic convex hull problem maintains $\operatorname{conv}(P(t))$ for
points moving along known trajectories. Its certificates are orientation
tests of hull edges and containment tests of interior points; its events are
the moments when a point enters or leaves the hull; and its repair
recomputes the affected tangent chains. The KDS of Basch, Guibas, and
Hershberger \cite{basch1999} is compact, local, responsive, and efficient.

\subsection{2. Definitions}

\begin{definition}[Kinetic Convex Hull]\label{def:kinetic_convex_hull}
  A \textbf{kinetic convex hull} structure maintains $\operatorname{conv}(P)$
  of $n$ moving points. Its certificates are of two kinds: for every hull
  edge $uv$, an orientation certificate asserting that all other points lie
  on the same side of the line through $u$ and $v$; and, for every interior
  point, a containment certificate asserting that it lies inside the hull.
  The \textbf{external events} are the times at which a point enters or
  leaves the hull.
\end{definition}

\subsection{3. Theory}

The certificates are sufficient: while every hull edge is a supporting line
and every interior point is contained, the stored hull is exact. A failure
needs one of two repairs --- a hull vertex moves so a neighbor leaves the
boundary (find a new tangent pair), or an interior point crosses a hull edge
(find the two tangents from the point to the hull). Both are local and use
the tangent routines of Chapter~\ref{ch:convex_hulls} on the dynamic hull of
Section~\ref{sec:dynamic_convex_hulls}.

\begin{theorem}[Kinetic Hull Bounds]\label{thm:kinetic_hull_bounds}
  The convex hull of $n$ points moving with constant velocity along lines
  can change $\Theta(n^2)$ times in the worst case. The kinetic convex hull
  of Basch, Guibas, and Hershberger \cite{basch1999} is compact ($O(n)$
  certificates), local (each point in $O(\log n)$), responsive (each event
  in $O(\log^2 n)$), and efficient: the number of internal events is within
  a constant factor of the number of external hull events.
\end{theorem}

\begin{proof}[Proof sketch]
  The lower bound adapts the crossing construction of
  Theorem~\ref{thm:kinetic_sort_events} to the plane, making each of
  $\Theta(n^2)$ pairwise crossings change the hull boundary. For the upper
  bound, correctness holds until the first certificate failure; the repair
  touches only the vertices of the failed edge and its neighbors, and the
  new certificates' failure times are roots of constant-degree polynomials
  (Section~\ref{sec:dynamic_certificates}). Each external event is triggered
  by a distinct failure, and internal events are bounded by locality.
\end{proof}

\subsection{4. Pseudocode}

\begin{algorithm}[htbp]
\caption{Kinetic Convex Hull}
\label{alg:kinetic_convex_hull}
\begin{algorithmic}[1]
\Require $n$ points with polynomial trajectories; a time horizon $[0, T]$.
\Ensure The convex hull as a function of time.
\Function{KineticHull}{$P$}
  \State build $\operatorname{conv}(P(0))$; create edge and containment certificates
  \State schedule each certificate's failure time in a kinetic heap $Q$
  \While{$Q$ is not empty and $Q.\text{min}() \le T$}
    \State $t \gets Q.\text{extract-min}()$; verify the failure exactly
    \If{an interior point crossed a hull edge}
      \State recompute the two tangents from the point to the hull; insert it
    \ElsIf{a hull vertex moved out}
      \State remove the vertex, connect its neighbors; recompute the local tangents
    \EndIf
    \State replace the certificates of the changed edges; reschedule their failure times
  \EndWhile
\EndFunction
\end{algorithmic}
\end{algorithm}

\subsection{5. Parameter Selection}

\begin{itemize}
  \item \textbf{Interior maintenance}: interior points are tracked by
    containment certificates so no point can exit without an event; only the
    closest potential exit matters, so certificates fail in the right order.
  \item \textbf{Degenerate events}: simultaneous entry/exit and collinear
    triples are resolved by the ordering rules of
    Section~\ref{sec:dynamic_certificates}.
  \item \textbf{Exact predicates}: orientation and containment tests use the
    adaptive exact methods of Section~\ref{sec:robust_predicates}; an
    approximate test misorders the event sequence.
\end{itemize}

\subsection{6. Complexity}

\begin{itemize}
  \item \textbf{Certificates}: $O(n)$; locality $O(\log n)$.
  \item \textbf{Responsiveness}: $O(\log^2 n)$ per event.
  \item \textbf{Events}: $\Theta(n^2)$ worst case for linear motion, so the
    total time over the horizon is $O(n^2 \log^2 n)$ and proportional to the
    number of hull changes plus polylogarithmic overhead.
\end{itemize}

\subsection{7. Usages}

\begin{itemize}
  \item \textbf{Motion planning}: hulls of moving obstacles support the
    continuous collision checks of Chapter~\ref{ch:motion_planning}.
  \item \textbf{Bounding volumes}: the hull of a moving cluster is the
    tightest convex bounding volume and is maintained kinetically instead of
    resampled.
  \item \textbf{Peeling}: a kinetic hull yields kinetic hull layers of
    Section~\ref{sec:convex_hull_peeling}.
\end{itemize}

\subsection{8. Testing Methods and Edge Cases}

\begin{itemize}
  \item \textbf{Static oracle}: at random times, recompute the hull from the
    positions and compare vertex sets.
  \item \textbf{Event oracle}: verify no point enters or leaves between
    consecutive events by re-checking all certificates at sample times.
  \item \textbf{Tangent correctness}: after each event, verify every point
    lies on one side of every hull edge.
\end{itemize}

\subsection{9. References}

The kinetic convex hull is developed in \cite{basch1999} and surveyed in
\cite{guibas1998}; the dynamic-hull subroutine is the structure of
Section~\ref{sec:dynamic_convex_hulls} (\cite{overmars1981,overmars1983}).
Textbook discussions appear in \cite{berg2008,mehta2005}.

\section{Kinetic Closest Pairs}
\label{sec:dynamic_kinetic_closest}

\subsection{1. Overview}

The kinetic closest-pair problem maintains the pair of points with minimum
mutual distance among moving points. The key observation is that the closest
pair is always a Delaunay edge (Chapter~\ref{ch:proximity}), so the problem
reduces to maintaining the Delaunay triangulation kinetically and reporting
its shortest edge. The kinetic Delaunay structure uses in-circle
certificates and edge flips, as discussed in the Voronoi chapter
(Section~\ref{sec:kinetic_voronoi}, Chapter~\ref{ch:voronoi_diagrams}).

\subsection{2. Definitions}

\begin{definition}[Kinetic Closest Pair]\label{def:kinetic_closest_pair}
  A \textbf{kinetic closest-pair} structure maintains, for each time $t$,
  the pair of distinct points minimizing the Euclidean distance, together
  with the minimum value. Its certificates are those of an underlying kinetic
  Delaunay triangulation (in-circle tests) plus the tournament certificates
  of Section~\ref{sec:dynamic_kds} used to select the shortest Delaunay
  edge.
\end{definition}

\subsection{3. Theory}

The reduction is exact: the closest pair is a Delaunay edge, and its length
is the minimum over all Delaunay edge lengths. The kinetic structure
maintains the Delaunay triangulation --- each edge carries an in-circle
certificate asserting that the circumcircles of its adjacent triangles are
empty --- and a kinetic tournament over the edge lengths. When an in-circle
certificate fails, the two affected triangles are flipped; when the
tournament's minimum changes, the closest pair changes.

\begin{theorem}[Kinetic Delaunay Complexity]\label{thm:kinetic_delaunay_bounds}
  The Delaunay triangulation (equivalently the Voronoi diagram) of $n$
  points moving with constant velocity along lines can change $\Theta(n^2)$
  times in the worst case \cite{albers1998}. For uniform linear motion, the
  number of events of a kinetic Delaunay triangulation is near-quadratic
  \cite{rubin2016}. Each event is an edge flip triggered by an in-circle
  certificate and repaired in $O(\log n)$ time.
\end{theorem}

\begin{theorem}[Kinetic Closest Pair Bounds]\label{thm:kinetic_closest_pair_bounds}
  The closest pair of $n$ points moving with constant velocity can change
  $\Omega(n^2)$ times in the worst case \cite{basch1997proximity}. The
  Delaunay-plus-tournament structure is compact ($O(n)$ certificates), local
  (each point in $O(\log n)$), and responsive ($O(\log^2 n)$ per event).
\end{theorem}

\begin{proof}[Proof sketch]
  The lower bound of \cite{basch1997proximity} adapts the crossing
  construction of Theorem~\ref{thm:kinetic_sort_events} so that the closest
  pair shifts among $\Theta(n^2)$ distinct pairs. For the upper bound, the
  triangulation is exact while all in-circle certificates hold (the
  empty-circle property of Chapter~\ref{ch:delaunay_triangulations}); a
  failure flips one edge and reschedules the $O(1)$ affected certificates,
  and the tournament adds $O(\log n)$ per minimum change. Each certificate
  is a degree-at-most-four polynomial under linear motion
  (Section~\ref{sec:dynamic_certificates}).
\end{proof}

\subsection{4. Pseudocode}

\begin{algorithm}[htbp]
\caption{Kinetic Closest Pair via Kinetic Delaunay}
\label{alg:kinetic_closest_pair}
\begin{algorithmic}[1]
\Require $n$ points with polynomial trajectories; time horizon $[0, T]$.
\Ensure The closest pair at all times in $[0, T]$.
\Function{KineticClosestPair}{$P$}
  \State build the Delaunay triangulation $DT$ of $P(0)$
  \For{each edge $e$ of $DT$}
    \State create an in-circle certificate for $e$; schedule its failure time
  \EndFor
  \State build a kinetic tournament $M$ over the edge lengths of $DT$
  \While{time $\le T$ and the event queue is nonempty}
    \State process the earliest event at time $t$
    \If{an in-circle certificate fails on edge $e$}
      \State flip $e$; remove the two old triangles, add the two new ones
      \State create certificates for the new edges; update $M$ on the changed lengths
    \EndIf
    \State report the minimum-length edge of $M$ as the closest pair at $t$
  \EndWhile
\EndFunction
\end{algorithmic}
\end{algorithm}

\subsection{5. Parameter Selection}

\begin{itemize}
  \item \textbf{In-circle polynomial}: under linear motion the certificate is
    a degree-four polynomial in $t$; with approximate failure times,
    re-prediction (Section~\ref{sec:dynamic_certificates}) keeps the
    structure exact.
  \item \textbf{All-nearest-neighbors}: the same Delaunay structure also
    supports kinetic nearest-neighbor queries per vertex; the tournament is
    replaced by a per-vertex report.
  \item \textbf{Higher dimensions}: in $\mathbb{R}^3$ the closest pair is
    not necessarily a Delaunay edge, so other certificate sets are used
    \cite{basch1997proximity}.
\end{itemize}

\subsection{6. Complexity}

\begin{itemize}
  \item \textbf{Certificates}: $O(n)$ (one per Delaunay edge).
  \item \textbf{Responsiveness}: $O(\log^2 n)$ per event.
  \item \textbf{Events}: $\Theta(n^2)$ worst case for linear motion
    \cite{albers1998}; near-quadratic for uniform linear motion
    \cite{rubin2016}.
\end{itemize}

\subsection{7. Usages}

\begin{itemize}
  \item \textbf{Particle simulation}: the closest approaching pair, with its
    distance used as a time-step safety bound.
  \item \textbf{Tracking}: the pair of moving agents with minimal separation,
    reported as a continuous function of time.
  \item \textbf{Collision detection}: a threshold certificate on a pair's
    distance fails exactly when a collision becomes possible.
\end{itemize}

\subsection{8. Testing Methods and Edge Cases}

\begin{itemize}
  \item \textbf{Brute-force oracle}: at dense sample times, compute the
    closest pair by all-pairs scanning (Section~\ref{sec:closest_pair},
    Chapter~\ref{ch:proximity}).
  \item \textbf{Delaunay invariant}: after every flip, verify the
    empty-circle property by exact in-circle predicates.
  \item \textbf{Coincident points and ties}: two points meeting produce a
    zero-distance closest pair (no division by zero in the failure-time
    computation); equal minima are resolved by a fixed rule.
\end{itemize}

\subsection{9. References}

The kinetic closest-pair structure and its lower bound are due to
\cite{basch1997proximity}; the Delaunay complexity bounds are
\cite{albers1998,rubin2016}; the framework is \cite{basch1999,guibas1998}.
Kinetic Voronoi diagrams are also treated in Section~\ref{sec:kinetic_voronoi}
of Chapter~\ref{ch:voronoi_diagrams}.

\section{Applications to Tracking and Simulation}
\label{sec:dynamic_applications}

The dynamic and kinetic structures of this chapter are the algorithmic core
of applications in which the world changes with time.

\begin{description}
  \item[Tracking] A moving target's position is sampled by sensors, and the
    system answers proximity queries about the current configuration.
    Dynamic range searching (Section~\ref{sec:rangesearch_dynamic},
    Chapter~\ref{ch:range_searching}) and the dynamic nearest-neighbor
    structures of Section~\ref{sec:dynamic_proximity} answer the queries;
    when trajectories are known, kinetic versions \emph{predict} the queries
    before they are asked.
  \item[Collision detection] In robotics and animation, moving objects must
    be checked for collisions, and the checks should run only when the
    configuration changes. The kinetic free-space approach of Agarwal,
    Basch, Guibas, Hershberger, and Zhang \cite{agarwal2002kinetic}
    maintains a tiling of the free space whose cells' validity is certified;
    a certificate failure signals a potential contact and triggers a local
    repair. This is the continuous-motion companion of the static checks of
    Chapter~\ref{ch:motion_planning}.
  \item[Simulation] Particle and agent-based models advance in discrete
    steps, but their combinatorial structure --- neighbors, hulls, closest
    pairs --- is best maintained kinetically so the time step is chosen
    adaptively from predicted events. The kinetic sorted list of
    Section~\ref{sec:dynamic_kinetic_sorting} is the canonical scheduling
    device for such simulations.
  \item[Geographic information systems] Moving objects (vehicles, animals)
    are queried over both space and time; kinetic range trees and the
    proximity queries of Section~\ref{sec:dynamic_kinetic_closest} extend
    the static GIS pipeline of Chapter~\ref{ch:applications_gis}.
\end{description}

Two lessons recur. First, the event-driven discipline of
Section~\ref{sec:dynamic_certificates} is only as good as its exact
predicates: a certificate that fails silently, or an event scheduled
slightly too late, corrupts every later answer, so production systems couple
the KDS with the adaptive exact predicates of
Section~\ref{sec:robust_predicates} and with re-prediction. Second, the
dynamic and kinetic views reinforce each other: a kinetic structure typically
reduces to a dynamic structure (the hull of Section~\ref{sec:dynamic_convex_hulls}
inside the kinetic hull of Section~\ref{sec:dynamic_kinetic_hulls}), so
progress on the dynamic problem directly improves the kinetic one. Open
problems --- robustness, event bounds for general motions, approximate
kinetics --- are surveyed in Chapter~\ref{ch:research_frontiers}.

\index{kinetic collision detection}\index{tracking}

\section{Exercises}
\label{sec:dynamic_exercises}

\begin{enumerate}
  \item Prove that the logarithmic method (Section~\ref{sec:dynamic_point_sets})
    converts a static structure with query time $Q(n)$ and construction time
    $C(n)$ into a semidynamic structure with amortized update cost
    $O(C(n)\log n / n)$ and query cost $O(Q(n)\log n)$, and derive the bounds
    for 2D orthogonal range counting.
  \item Argue that the Overmars--van Leeuwen dynamic hull
    (Section~\ref{sec:dynamic_convex_hulls}) answers a membership query in
    $O(\log n)$ time by binary search on the hull tree, and give an input on
    which a naive point-in-polygon test takes $O(n)$.
  \item Show that a direction query in the dynamic hull is answered by a
    single root-to-leaf walk using stored bridges, and explain why a node's
    bridge must be recomputed after any update in its subtree.
  \item Construct the two-velocity family of Theorem~\ref{thm:kinetic_sort_events}
    explicitly for $n = 6$ and list the times at which the sorted order
    changes; verify that the kinetic sorted list processes exactly those
    events.
  \item Prove the four quality measures (Definition~\ref{def:kinetic_properties})
    for the kinetic sorted list: compactness, locality $O(1)$, responsiveness
    $O(\log n)$ per event, and efficiency $O(1)$ relative to order changes.
  \item Verify that the in-circle certificate under linear motion is a
    polynomial of degree at most four by writing the certificate determinant
    and counting degrees, and state the degree of an orientation certificate.
  \item Design a kinetic structure maintaining the pair of points with
    maximum distance (the kinetic diameter) using a kinetic tournament, and
    discuss whether its event count can be quadratic in the worst case.
  \item Modify the kinetic closest-pair algorithm
    (Section~\ref{sec:dynamic_kinetic_closest}) to handle two points that
    coincide at a time, and state what the exact predicates and the
    failure-time computation must return.
\end{enumerate}
